\section{引言：这一讲到底想解决什么}

上一讲已经把生成式建模的前半张图铺好了：autoregressive 模型直接分解 $p(x)$，VAE 用 latent variable 和 ELBO 近似建模数据分布，GAN 走向隐式密度，diffusion 则把生成变成逐步去噪。今天这一讲不是“再加一个模型名字”，而是把这张图的后半边补完整，并且把现代图像、视频生成系统为什么会变成“模块拼装”的样子讲清楚。

\begin{figure}[H]
\centering
\includegraphics[width=0.95\textwidth]{slides-images/slide-000.jpg}
\caption{标题页：本讲延续上一讲的生成式建模主线，重点放在 GAN、diffusion、latent diffusion 和现代多模态生成管线\protect\footnotemark}
\end{figure}
\footnotetext{Slides 第1页。}

\begin{knowledgebox}{本讲主线}
这讲的核心不是“哪个模型更酷”，而是把同一个问题拆成四个层次：
\begin{itemize}
    \item \textbf{生成机制}：noise 如何变成 data
    \item \textbf{训练机制}：监督信号从哪里来
    \item \textbf{采样机制}：如何在推理时真正生成样本
    \item \textbf{工程机制}：为什么现代系统几乎一定要拼接 VAE、GAN、diffusion、Transformer 和 text encoder
\end{itemize}
\end{knowledgebox}

\subsection{先把生成模型的家族图再看一遍}

从讲义角度，最容易混淆的是“显式密度”和“隐式密度”这两个维度。显式密度模型试图直接写出或近似写出 $p(x)$；隐式密度模型不强调概率值本身，而强调能否从模型里采样出看起来像真的样本。前者重 likelihood，后者重 sample quality。

\begin{figure}[H]
\centering
\includegraphics[width=0.95\textwidth]{slides-images/slide-003.jpg}
\caption{上一讲的回顾：判别式模型学 $p(y\mid x)$，生成式模型学 $p(x)$ 或 $p(x\mid y)$\protect\footnotemark}
\end{figure}
\footnotetext{Slides 第3页。}

\begin{figure}[H]
\centering
\includegraphics[width=0.95\textwidth]{slides-images/slide-004.jpg}
\caption{生成模型家族树：autoregressive 和 VAE 属于显式密度，GAN 和 diffusion 则分别代表隐式采样和逐步去噪\protect\footnotemark}
\end{figure}
\footnotetext{Slides 第4页。}

\begin{table}[H]
\centering
\small
\renewcommand{\arraystretch}{1.15}
\resizebox{\textwidth}{!}{%
\begin{tabular}{p{2.9cm}p{2.5cm}p{3.2cm}p{2.6cm}p{4.2cm}}
\toprule
\textbf{方法} & \textbf{密度视角} & \textbf{训练信号} & \textbf{采样方式} & \textbf{最典型的工程特征} \\
\midrule
Autoregressive & 显式密度 & 最大化对数似然 & 逐 token 采样 & 似然好看，但 raw pixels 上极慢 \\
VAE & 近似显式密度 & ELBO & 一步解码 & latent 结构清楚，但重建容易糊 \\
GAN & 隐式密度 & 对抗博弈 & 一步生成 & 样本锐利，但训练极不稳定 \\
Diffusion / flow & 逐步运输 / 去噪 & 回归向量场或 score & 多步积分 & 训练稳，但采样慢 \\
Latent diffusion & 压缩后的 diffusion & 先学 autoencoder，再学 diffusion & latent 上多步采样 & 这是现代图像/视频系统的主干 \\
\bottomrule
\end{tabular}%
}
\caption{这一讲的知识地图：不同生成模型本质上都在解决“如何从简单噪声到复杂数据”的同一个问题，只是训练信号和采样代价不同}
\end{table}

\subsection{为什么今天的主角是 GAN 和 diffusion}

autoregressive 和 VAE 已经在上一讲讲过，这一讲重点是两条后半路：GAN 让我们看到了“不要显式建模密度也能生成很漂亮的样本”，diffusion 则让我们看到“把生成拆成很多个小步骤，训练会稳很多”。这两条路最后没有互相消灭，反而在工业里合流成一套组合拳。

\begin{figure}[H]
\centering
\includegraphics[width=0.95\textwidth]{slides-images/slide-006.jpg}
\caption{今天的目标页：本讲从上一讲的三类模型继续往下走，进入 GAN 和 diffusion\protect\footnotemark}
\end{figure}
\footnotetext{Slides 第7页。}

\begin{importantbox}{一眼看懂本讲的结论}
如果只记一件事，就是这句：\textbf{现代生成式系统不是“选一个模型”，而是“组合一组互补模块”}。GAN 提供锐利感，VAE 提供压缩与 latent，diffusion 提供稳定训练，Transformer 提供可扩展主干，text encoder 提供条件控制。
\end{importantbox}

